\documentclass[10pt,a4paper]{amsart}
\usepackage[margin=19mm]{geometry}
\usepackage{amsmath,amssymb,mathtools}
\usepackage[scr=rsfs,cal=euler]{mathalfa}
\usepackage{xfrac}
\usepackage[unicode,hidelinks,bookmarks=false]{hyperref}
\newcommand{\ie}{i.e.\ }
\newcommand{\cf}{cf.\ }
\newcommand{\resp}{respectively\ }
\newcommand{\Subsection}{\subsection}
\providecommand{\nobreakdash}{\nobreak\hbox{-}}
\setlength{\emergencystretch}{2em}
\multlinegap0pt
\usepackage{amssymb}
\usepackage{xcolor}
\usepackage{algorithm}\usepackage{algpseudocode}
\newtheorem{theorem}{Theorem}
\newtheorem{lemma}{Lemma}
\def\bW{\mathbf{W}}
\def\bX{\mathbf{X}}
\def\bb{\mathbf{b}}
\def\bx{\mathbf{x}}
\def\by{\mathbf{y}}
\def\bz{\mathbf{z}}
\def\intd{\mathrm{d}}
\title{High-Dimensional Enhanced Sampling via Regularized Path-Dependent McKean--Vlasov Dynamics using Tensor Density Approximation}
\author{Liyao Lyu, Siyu Guo, Huan Lei}
\date{Source: arXiv:2605.03080v1 (CC BY 4.0)}
\begin{document}
\maketitle

\noindent\emph{Source and licence.} High-Dimensional Enhanced Sampling via Regularized Path-Dependent McKean--Vlasov Dynamics using Tensor Density Approximation. Version arXiv:2605.03080v1 (CC BY 4.0), distributed by arXiv under the Creative Commons Attribution 4.0 International licence, http://creativecommons.org/licenses/by/4.0/, as recorded in the arXiv licence metadata for this exact version and shown on the abstract page https://arxiv.org/abs/2605.03080v1. Full licence notice: \texttt{licenses/arxiv-2605.03080-CC-BY-4.0.txt}.\par\smallskip

\noindent\textit{Editorial scope.} Counted unit: the article's lemma lem:lip\_cont (source0.tex lines 383--405). The article's complete original proof of it is delivered with it (source0.tex lines 407--476, one \texttt{proof} environment of 70 lines). No mathematics is written, repaired or re-derived: the statement and the proof are the article's own lines, in the article's own notation, and no retained line is substituted or re-flowed.  Retained uncounted as the source-local context the proof needs: lines 360--362; lines 375--378.  Nothing was omitted from any retained span, so the package carries no omission record.  No retained line contains a citation, so the wrapper carries no bibliography and no bibliography entry.  No retained line contains a figure environment, an image-inclusion command or drawing code, so no image is delivered and the package declares no asset dependency.  Editorial formatting only, in addition: an AMS \texttt{amsart} preamble that restates the article's own declarations and macros used by the retained lines, namely the environments \texttt{theorem}, \texttt{lemma} on separate counters, together with the standard packages those lines need.  The retained environments print in this document as Theorem 1, Lemma 1 in order of appearance, whereas the article prints the same items with the numbering of its own full document.  The complete native source of the article as distributed in the arXiv e-print for this exact version is bundled unchanged as \texttt{sources/199768/source0.tex}.
\begin{equation}\label{equ:regularized_MV_x}
    \intd \bX_t = -\frac{1}{\gamma}\nabla\!\left(U(\bX_t) + \frac{\alpha}{\beta}\log K_{\tau,\epsilon}\!\bigl((\varphi^\delta * \mu_{\xi,t})(\xi(\bX_t))\bigr)\right) \intd t + \sqrt{\frac{2}{\beta\gamma}}\,\intd \bW_t.
\end{equation}

\begin{theorem}\label{thm:mk_well}
Assume that $\nabla U(\bx)$ is globally Lipschitz, that $\xi:\mathbb R^{d}\to\mathbb R^m$ belongs to $C^2$ with bounded $\nabla \xi$ and globally Lipschitz $\nabla^2 \xi$, and that $\varphi^\delta\in C^2(\mathbb R^m)$ with bounded derivatives. 
Then, for any deterministic initial condition \(\bX_0=\bx\), the McKean--Vlasov SDE~\eqref{equ:regularized_MV_x} admits a unique strong solution.
\end{theorem}

\begin{lemma} \label{lem:lip_cont}
Under the assumptions of Theorem~\ref{thm:mk_well},
let 
\begin{equation}
\bb(\bX,\mu)
=
-\frac{1}{\gamma}\nabla \left( U(\bX)
+
\frac{\alpha}{\beta}
\log\!
K_{\tau,\epsilon}\left((\varphi^\delta * \mu_{\xi})(\xi(\bX))\right)
\right), \quad \sigma = \sqrt{\frac{2}{\beta\gamma}}.
\label{eq:b_drift}
\end{equation}
Then there exists $C>0$ such that
\begin{equation}\label{equ:lip}
   \|\bb(\bX_1,\mu_1)-\bb(\bX_2,\mu_2)\|
\le
C\Big(\|\bX_1-\bX_2\|+W_2(\mu_1,\mu_2)\Big) 
\end{equation}
for any $\bX_1,\bX_2\in \mathbb R^{d}$ and $\mu_1,\mu_2\in \mathcal P_2(\mathbb{R}^{d})$.
As a consequence, $\|\bb(\bX,\mu)\| \leq C(1+\|\bX\|)$.
\end{lemma}

\begin{proof}
We decompose the drift into a potential and a bias component. Since $\nabla U$ is globally Lipschitz,
\[
\left\|-\frac1\gamma \nabla U(\bX_1)+\frac1\gamma \nabla U(\bX_2)\right\|
\le
\frac{\mathrm{Lip}(\nabla U)}{\gamma}\|\bX_1-\bX_2\|.
\]
It remains to estimate the bias term. Define
\[
f_\mu(\bz):=(\varphi^\delta * \mu_{\xi})(\bz)
=
\int_{\mathbb R^{d}} \varphi^\delta(\bz-\xi(\by))\,\mu(\intd \by),
\quad
H_\mu(\bz):=\nabla_\bz \log\!\big(K_{\tau,\epsilon}(f_\mu(\bz))\big).
\]
For any $\bz\in\mathbb R^m$ and any coupling $\pi\in\Pi(\mu_1,\mu_2)$,
\[
f_{\mu_1}(\bz)-f_{\mu_2}(\bz)
=
\iint
\Big(
\varphi^\delta(\bz-\xi(\by_1))
-
\varphi^\delta(\bz-\xi(\by_2))
\Big)\,\pi(\intd \by_1,\intd \by_2).
\]
Since $\nabla \varphi^\delta$ is bounded and $\xi$ is Lipschitz (because $\nabla \xi$ is bounded), taking the infimum over $\pi$ and supremum over $\bz$ yields
\[
\|f_{\mu_1}-f_{\mu_2}\|_{L^\infty(\mathbb R^m)}
\le
C\,W_1(\mu_1,\mu_2)
\le
C\,W_2(\mu_1,\mu_2).
\]
Similarly, since $\nabla \varphi^\delta$ is $C^1$ with bounded derivative,
$\|\nabla f_{\mu_1}-\nabla f_{\mu_2}\|_{L^\infty(\mathbb R^m)}
\le
C\,W_2(\mu_1,\mu_2)$.

Now set $h(r):=\log K_{\tau,\epsilon}(r)$.
Because $K_{\tau,\epsilon}\ge \epsilon>0$, both $h'$ and $h''$ are bounded. Therefore
\[
\|H_{\mu_1}-H_{\mu_2}\|_{L^\infty(\mathbb R^m)}
\le
C\Big(
\|f_{\mu_1}-f_{\mu_2}\|_{L^\infty}
+
\|\nabla f_{\mu_1}-\nabla f_{\mu_2}\|_{L^\infty}
\Big)
\le
C\,W_2(\mu_1,\mu_2).
\]

Next we control the dependence on $\bX$. Since $\varphi^\delta\in C^2$ with bounded derivatives, $H_\mu$ is globally Lipschitz in $\bz$, uniformly in $\mu$.
Using the boundedness of $\nabla \xi$ and the Lipschitz continuity of $\nabla \xi$ (which follows from the boundedness of $\nabla^2 \xi$), we obtain
\begin{align*}
&\|\nabla \xi(\bX_1)^\top H_{\mu_1}(\xi(\bX_1))
-
\nabla \xi(\bX_2)^\top H_{\mu_2}(\xi(\bX_2))\| \\
&\qquad \le
C\|\bX_1-\bX_2\|
+
C\|H_{\mu_1}(\xi(\bX_1))-H_{\mu_1}(\xi(\bX_2))\|
+
C\|H_{\mu_1}(\xi(\bX_2))-H_{\mu_2}(\xi(\bX_2))\| \\
&\qquad \le
C\|\bX_1-\bX_2\|+C\,W_2(\mu_1,\mu_2).
\end{align*}
Combining with the Lipschitz estimate for $\nabla U$ gives the result.
\end{proof}

\end{document}
